\clearpage\SourceLeaf{172}\pdfbookmark[0]{Neue Methode zur näherungsweisen Bestimmung der Werte von Integralen}{seventh-paper}\thispagestyle{empty}
\begin{center}
NEUE METHODE

\medskip
ZUR NÄHERUNGSWEISEN BESTIMMUNG

DER WERTE VON INTEGRALEN

VON

CARL FRIEDRICH GAUSS

DER KÖNIGLICHEN GESELLSCHAFT DER WISSENSCHAFTEN AM 16. SEPTEMBER 1814 VORGELEGT.

\par\bigskip\rule{.7\linewidth}{.4pt}\par\medskip
Commentationes societatis regiae scientiarum Gottingensis recentiores. Vol.~\textsc{iii}.

Gottingae \textsc{mdcccxvi}.
\par\medskip\rule{.7\linewidth}{.4pt}
\end{center}
\clearpage\fancyhead[L]{\small NÄHERUNGSWERTE VON INTEGRALEN}
\SourceLeaf{174}
\begin{center}
NEUE METHODE

ZUR NÄHERUNGSWEISEN BESTIMMUNG

DER WERTE VON INTEGRALEN.

\medskip\rule{.1\linewidth}{.4pt}
\end{center}
\SeventhArticle{1}
Unter den zur näherungsweisen numerischen Bestimmung von Integralen vorgeschlagenen Methoden nehmen die Regeln einen hervorragenden Platz ein, die \textsc{Cotes} im Anschluss an den großen \textsc{Newton} entwickelt hat. Wird nämlich der Wert des Integrals $\int y\,\mathrm{d}x$ von $x=g$ bis $x=h$ gesucht, so sind die Werte von $y$ für diese beiden äußersten Werte von $x$ und für beliebig viele Zwischenwerte, die vom ersten bis zum letzten in gleichen Schritten fortschreiten, mit gewissen numerischen Koeffizienten zu multiplizieren. Die Summe der Produkte, mit $h-g$ multipliziert, liefert dann das gesuchte Integral, und zwar umso genauer, je mehr Glieder bei dieser Rechnung verwendet werden. Da die Grundsätze dieser Methode, die von den Geometern anscheinend seltener angewendet wird, als sie es verdient, meines Wissens nirgends ausführlicher erläutert worden sind, wird es unserem Vorhaben entsprechen, einige Bemerkungen darüber vorauszuschicken.

\SeventhArticle{2}
Die Anzahl der Glieder, die verwendet werden sollen, sei $n+1$, und wir setzen $h-g=\Delta$, sodass die Werte von $x$ gleich $g$, $g+\frac{\Delta}{n}$, $g+\frac{2\Delta}{n}$, $g+\frac{3\Delta}{n}$ usw. bis $g+\Delta$ sind, wobei die entsprechenden Werte von $y$ jeweils $A$, $A'$, $A''$, $A'''$ usw. bis $A^{(n)}$ sind. Schließlich werde allgemein $x=g+\Delta t$ gesetzt, sodass $y$ auch als Funktion von $t$ betrachtet werden kann. Wir bezeichnen mit $Y$ die folgende Funktion
\SourceLeaf{175}
\[
\begin{aligned}
&A.\;\frac{(nt-1)(nt-2)(nt-3)\ldotp\ldotp\ldotp\ldotp(nt-n)}{(-1)\quad(-2)\, .\,(-3)\ldotp\ldotp\ldotp\ldotp(-n)}\\
{}+{}&A'.\;\frac{nt\, .\,(nt-2)\, .\,(nt-3)\ldotp\ldotp\ldotp\ldotp\ldotp\ldotp(nt-n)}{1\, .\,(-1)\quad(-2)\ldotp\ldotp\ldotp\ldotp\ldotp\ldotp(1-n)}\\
{}+{}&A''.\;\frac{nt\, .\,(nt-1)\, .\,(nt-3)\ldotp\ldotp\ldotp\ldotp\ldotp\ldotp(nt-n)}{2.\quad1.\quad(-1)\ldotp\ldotp\ldotp\ldotp\ldotp\ldotp(2-n)}\\
{}+{}&A'''.\;\frac{nt\, .\,(nt-1)\, .\,(nt-2)\ldotp\ldotp\ldotp\ldotp\ldotp\ldotp(nt-n)}{3.\quad2.\quad1\ldotp\ldotp\ldotp\ldotp\ldotp\ldotp(3-n)}\\
{}+{}&\text{usw.}\\
{}+{}&A^{(n)}.\;\frac{nt(nt-1)\, .\,(nt-2)\ldotp\ldotp\ldotp\ldotp\ldotp\ldotp(nt-n+1)}{n\, .\,(n-1)\quad(n-2)\ldotp\ldotp\ldotp\ldotp\ldotp\ldotp1}
\end{aligned}
\]
oder $\sum\dfrac{A^{(\mu)}T^{(\mu)}}{M^{(\mu)}}$, wobei $\mu$ die einzelnen ganzen Zahlen $0,1,2,3\ldotp\ldotp\ldotp n$ durchläuft,
\[
T^{(\mu)}=\frac{nt(nt-1)(nt-2)(nt-3)\ldotp\ldotp\ldotp\ldotp(nt-n)}{nt-\mu}
\]
\[
M^{(\mu)}\quad\text{Wert von }T\text{ für }nt=\mu.
\]
Offenbar ist $Y$ ein Polynom in $t$ vom Grad $n$, und seine Werte für die einzelnen $n+1$ Werte von $t$, nämlich $0,\frac1n,\frac2n,\frac3n\ldotp\ldotp\ldotp\ldotp1$, sind gleich den Werten von $y$. Ferner ist klar: Ist $Y'$ ein anderes Polynom, das für dieselben Werte mit $y$ übereinstimmt, so verschwindet $Y'-Y$ für diese Werte und ist daher durch die Faktoren $t,t-\frac1n,t-\frac2n,t-\frac3n\ldotp\ldotp\ldotp\ldotp t-1$ und folglich auch durch deren Produkt (das den Grad $n+1$ hat) teilbar. Somit muss $Y'$, sofern es nicht völlig mit $Y$ identisch ist, gewiss einen höheren Grad haben; mit anderen Worten ist $Y$ unter allen Polynomen, deren Grad $n$ nicht überschreitet, das einzige, das für jene $n+1$ Werte mit $y$ übereinstimmt. Bricht also $y$, in eine Reihe nach Potenzen von $t$ entwickelt, vollständig vor dem Glied ab, das $t^{n+1}$ enthält, so ist diese Funktion mit $Y$ identisch. Konvergiert die Reihe dagegen wenigstens so schnell, dass die folgenden Glieder vernachlässigt werden dürfen, so kann die Funktion $Y$ zwischen den Grenzen $t=0,t=1$ oder $x=g,x=h$ an die Stelle von $y$ treten.

